\section*{第 \ref{ch:b1:faults-folds-fabrics} 章——断层、褶皱与组构}

\begin{solution}{exo:b1:faults-folds-fabrics:1}
(a) \omterm{def:b1:faults-folds-fabrics:fault-types}{正断层}。(b) \omterm{def:b1:faults-folds-fabrics:fault-types}{逆断层}，近乎纯倾滑。(c) \omterm{def:b1:faults-folds-fabrics:fault-types}{左行}
\omterm{def:b1:faults-folds-fabrics:fault-types}{走滑断层}。(d) \omterm{def:b1:faults-folds-fabrics:fault-types}{右行}\omterm{def:b1:faults-folds-fabrics:fault-types}{走滑断层}。(e) \omterm{def:b1:faults-folds-fabrics:fault-types}{斜滑断层}：
\omterm{def:b1:faults-folds-fabrics:fault-elements}{上盘}向下并沿\omterm{def:b1:geological-maps:attitude}{走向}向左移动；其走滑
分量占\omterm{def:b1:faults-folds-fabrics:fault-elements}{总滑距}的 $\cos\ang{40} = 0.77$，大于其
倾滑分量 $0.64$，因此称为带正断分量的\omterm{def:b1:faults-folds-fabrics:fault-types}{左行}断层。
(f) \omterm{def:b1:geological-maps:attitude}{倾角}为 \ang{25}，小于 \ang{45}，该\omterm{def:b1:faults-folds-fabrics:fault-types}{逆断层}为
\omterm{def:b1:faults-folds-fabrics:fault-types}{逆冲断层}。
\end{solution}

\begin{solution}{exo:b1:faults-folds-fabrics:2}
当 $|\sin r| = 1$ 时，\omterm{def:b1:cross-sections:throw}{垂直断距}为 $300\sin\ang{55} = \qty{245.7}{m}$，且
\omterm{def:b1:cross-sections:throw}{水平断距}为 $300\cos\ang{55} = \qty{172.1}{m}$。
\end{solution}

\begin{solution}{exo:b1:faults-folds-fabrics:3}
对于\omterm{def:b1:faults-folds-fabrics:fold-classes}{直立褶皱}，其两翼分别以 $\delta_1$ 和 $\delta_2$ 朝相反
方向倾斜，$\iota = \ang{180} - \delta_1 - \delta_2$。(a) \ang{110}：开阔。
(b) \ang{35}：紧闭。(c) \ang{7}：等斜。
\end{solution}

\begin{solution}{exo:b1:faults-folds-fabrics:4}
\omterm{def:b1:faults-folds-fabrics:fault-types}{正断层}\omterm{def:b1:geological-maps:attitude}{倾角}为 $\ang{45} + \ang{35}/2 = \ang{62.5}$，\omterm{def:b1:faults-folds-fabrics:fault-types}{逆冲断层}
\omterm{def:b1:geological-maps:attitude}{倾角}为 $\ang{45} - \ang{17.5} = \ang{27.5}$。共轭\omterm{def:b1:faults-folds-fabrics:fault-types}{走滑断层}各以
\ang{27.5} 的角度与 $\sigma_1$ 相交，分居其两侧，二者的夹角为 \ang{55}。
\end{solution}

\begin{solution}{exo:b1:faults-folds-fabrics:5}
(a) 走滑分量为 $400\cos\ang{30} = \qty{346.4}{m}$，沿
\omterm{def:b1:geological-maps:attitude}{走向}；倾滑分量为 $400\sin\ang{-30} = \qty{-200}{m}$，即
沿\omterm{def:b1:geological-maps:attitude}{倾向}向下 \qty{200}{m}；\omterm{def:b1:cross-sections:throw}{垂直断距}为 $200\sin\ang{50} = \qty{153.2}{m}$，且
\omterm{def:b1:cross-sections:throw}{水平断距}为 $200\cos\ang{50} = \qty{128.6}{m}$。(b) \omterm{def:b1:geological-maps:attitude}{走向}为 $000$，故
$\mathbf s = (0, 1, 0)$; $\mathbf u = (-\cos\ang{50}, 0, \sin\ang{50}) =
(-0.643, 0, 0.766)$。则
$\overrightarrow{AA'} = 400(0.866\,\mathbf s - 0.5\,\mathbf u) =
(128.6, 346.4, -153.2)$（单位为米）。其\omterm{def:b1:field-lab-methods-1:trend-plunge}{倾伏向}为 $\arctan(128.6/346.4) =
\ang{20.4}$，\omterm{def:b1:field-lab-methods-1:trend-plunge}{倾伏角}为 $\arcsin(153.2/400) = \ang{22.5}$，
这也等于 $\arcsin(\sin\ang{30}\sin\ang{50})$。(c) \omterm{def:b1:faults-folds-fabrics:fault-types}{斜滑断层}，\omterm{def:b1:faults-folds-fabrics:fault-elements}{侧伏角}
介于 $0$ 与 \ang{-90} 之间，且更接近 $0$，为\omterm{def:b1:faults-folds-fabrics:fault-types}{左行}
断层，带正断分量。
\end{solution}

\begin{solution}{exo:b1:faults-folds-fabrics:6}
(a) 两断层彼此成 \ang{60}；$\sigma_1$ 平分
锐角，故其沿 $050$（及 $230$）水平；$\sigma_3$ 水平
沿 $140$ 方向；$\sigma_2$ 竖直，沿两竖直断层相交的
交线。各断层与 $\sigma_1$ 成 \ang{30}，故 $\ang{45} - \phi/2 =
\ang{30}$，且 $\phi = \ang{30}$。(b) 两断层将地面切割为四
个楔块；含 $\sigma_1$ 的楔块，朝向 $050$ 与 $230$，相向运动，
而含 $\sigma_3$ 的楔块则相背运动，朝向 $140$
与 $320$。取\omterm{def:b1:geological-maps:attitude}{走向}为 $020$ 的断层，位于朝向 $050$ 的楔块（其
东侧）与朝向 $320$ 的楔块（其西侧）之间：相对于西侧楔块，
东侧楔块朝向 $230$ 运动，并远离 $320$，二者均有
沿断层朝向 $200$ 的分量。站在西侧楔块上的观察者，
面朝东，看到东侧楔块朝 $200$ 运动，即向右，
故\omterm{def:b1:geological-maps:attitude}{走向} $020$ 的断层为\omterm{def:b1:faults-folds-fabrics:fault-types}{右行}。同理，或由关于
$\sigma_1$ 的对称性，可得\omterm{def:b1:geological-maps:attitude}{走向} $080$ 的断层为\omterm{def:b1:faults-folds-fabrics:fault-types}{左行}。
\end{solution}

\begin{solution}{exo:b1:faults-folds-fabrics:7}
向上的法线为
\begin{align*}
  \mathbf n_1 &= (\sin\ang{40}\sin\ang{120}, \sin\ang{40}\cos\ang{120}, \cos\ang{40}) = (0.557, -0.321, 0.766),\\
  \mathbf n_2 &= (\sin\ang{30}\sin\ang{260}, \sin\ang{30}\cos\ang{260}, \cos\ang{30}) = (-0.492, -0.087, 0.866).
\end{align*}
二者的叉积为
$(-0.211, -0.859, -0.207)$，指向下方，因此\omterm{def:b1:faults-folds-fabrics:fold-elements}{枢纽}线
\omterm{def:b1:field-lab-methods-1:trend-plunge}{倾伏角}为 $\arcsin 0.207/0.909 = \ang{13.1}$，\omterm{def:b1:field-lab-methods-1:trend-plunge}{倾伏向}为 $\ang{180} + \arctan(0.211/0.859) =
\ang{193.8}$。$\mathbf n_1\cdot\mathbf n_2 = 0.417$，两法线之间的
夹角为 \ang{65.3}，故 $\iota = \ang{114.7}$：\omterm{def:b1:faults-folds-fabrics:fold-classes}{开阔褶皱}。
\end{solution}

\begin{solution}{exo:b1:faults-folds-fabrics:8}
$\mathbf n_1 = (0.321, 0.383, 0.866)$ 为层理面的，而
$\mathbf n_2 = (0.337, -0.925, 0.174)$ 为\omterm{def:b1:faults-folds-fabrics:foliation}{劈理}面的。二者的叉积，
取向下者，得一直线，\omterm{def:b1:field-lab-methods-1:trend-plunge}{倾伏角}为 \ang{25.4}，\omterm{def:b1:field-lab-methods-1:trend-plunge}{倾伏向}为 \ang{074.8}。由
\cref{prop:b1:faults-folds-fabrics:intersection}，此即
\omterm{def:b1:faults-folds-fabrics:fold-elements}{枢纽}线的方向：\omterm{def:b1:geological-maps:fold}{褶皱}大致以 \ang{25} 向东北东倾伏，这
一结论可由两翼中任一其他\omterm{def:b1:geological-maps:geological-map}{露头}证实。
\end{solution}

\begin{solution}{exo:b1:faults-folds-fabrics:9}
(a) 站在断层以南、面朝北的观察者，看到岩墙北段
向东错开，即向右：该断层为\omterm{def:b1:faults-folds-fabrics:fault-types}{右行}。(b) 岩墙为
与断层垂直的竖直平面，故其沿\omterm{def:b1:geological-maps:attitude}{走向}的错距
即走滑分量 \qty{350}{m}，若
\omterm{def:b1:faults-folds-fabrics:fault-rocks}{擦痕}水平，则即为全部滑动量。(c) 此时滑动与\omterm{def:b1:geological-maps:attitude}{走向}成 \ang{20} 角，
位于竖直断层面内，其水平分量仍为 \qty{350}{m}，
故\omterm{def:b1:faults-folds-fabrics:fault-elements}{总滑距}为 $350/\cos\ang{20} = \qty{372.5}{m}$，其竖直
分量为 $350\tan\ang{20} = \qty{127.4}{m}$。北侧块体向东移动，
且\omterm{def:b1:faults-folds-fabrics:fault-rocks}{擦痕}朝东下降，故北侧亦向下移动，下移量为
\qty{127.4}{m}。竖直岩墙无法显示这种竖直运动，因为
沿其自身滑动，其迹线不变。
\end{solution}

\begin{solution}{exo:b1:faults-folds-fabrics:10}
(a) $\ang{90} - \theta$ 即断层面法线与 $\sigma_1$ 的夹角，故由
\cref{prop:b1:stress-strain-rheology:mohr}可得
$\sigma_n' = \sigma_1'\sin^2\theta + \sigma_3'\cos^2\theta$，且
$\tau = (\sigma_1' - \sigma_3')\sin\theta\cos\theta$。滑动开始于
$\tau = \mu\sigma_n'$ 时；两边除以 $\sigma_3'\cos^2\theta$，并记
$R = \sigma_1'/\sigma_3'$，$(R - 1)\tan\theta = \mu(R\tan^2\theta + 1)$，故
$R\tan\theta(1 - \mu\tan\theta) = \tan\theta + \mu$，即所求之式。
(b) 取 $\sigma_1$ 竖直，$\theta$ 为 \ang{90} 减去\omterm{def:b1:geological-maps:attitude}{倾角}：
$\theta = \ang{30}$ 时 $R = 3.12$；$\theta = \ang{45}$ 时 $R = 4.00$；
$\theta = \ang{55}$ 时 $R = 9.92$。(c) 分母为零的条件为
$\tan\theta = 1/\mu$，即 $\theta = \ang{59.0}$；超过此值，$R$ 将为负：
正\omterm{def:b1:stress-strain-rheology:lithostatic}{有效应力}之比不论如何都无法使断层滑动，即使
$\sigma_1'$ 增至多大。\omterm{def:b1:geological-maps:attitude}{倾角}小于 \ang{31.0} 的\omterm{def:b1:faults-folds-fabrics:fault-types}{正断层}无法
在竖直 $\sigma_1$、$\mu = 0.6$ 的情况下滑动。(d) 取向最有利的面
$\ang{45} - \phi/2$ 即其与 $\sigma_1$ 的夹角。取 $\mu = 0.6$，$\phi = \ang{31.0}$：
\omterm{def:b1:faults-folds-fabrics:fault-types}{正断层}\omterm{def:b1:geological-maps:attitude}{倾角} \ang{60.5}，\omterm{def:b1:faults-folds-fabrics:fault-types}{逆冲断层} \ang{29.5}。取 $\mu = 0.85$，
$\phi = \ang{40.4}$：\omterm{def:b1:faults-folds-fabrics:fault-types}{正断层}、\omterm{def:b1:faults-folds-fabrics:fault-types}{逆冲断层}的\omterm{def:b1:geological-maps:attitude}{倾角}分别为 \ang{65.2} 与 \ang{24.8}。
\end{solution}

\begin{solution}{exo:b1:faults-folds-fabrics:11}
朝北看，东在右侧，向西倾的面向右
抬升。(a) 在 B 处，层理与\omterm{def:b1:faults-folds-fabrics:foliation}{劈理}\omterm{def:b1:geological-maps:attitude}{倾向}相同，但更陡：
由\cref{prop:b1:faults-folds-fabrics:cleavage-bedding}(a)，B 位于
倒转翼上；A 与 C 处层理较缓，
为正立；(b) \omterm{def:b1:faults-folds-fabrics:foliation}{劈理}\omterm{def:b1:geological-maps:attitude}{倾角}为 \ang{60}，A 与 C 处层理\omterm{def:b1:geological-maps:attitude}{倾角}为
\ang{40}，两者之间锐角旋转为顺时针，
角度为 \ang{20}，使\omterm{def:b1:faults-folds-fabrics:foliation}{劈理}顶端移向东侧；B 处层理
\omterm{def:b1:geological-maps:attitude}{倾角}为 \ang{80}，旋转为逆时针，角度为 \ang{20}，使其
移向西侧。(c) 旋转方向在 A、B 之间及 B、C 之间各变化一次：
每段之间各有一个\omterm{def:b1:faults-folds-fabrics:fold-elements}{转折端}。A 处正立，其
\omterm{def:b1:faults-folds-fabrics:foliation}{劈理}顶端向东，指向翼顶的\omterm{def:b1:geological-maps:fold}{背斜}：A、B 之间的\omterm{def:b1:faults-folds-fabrics:fold-elements}{转折端}
为\omterm{def:b1:geological-maps:fold}{背斜}，B、C 之间的为\omterm{def:b1:geological-maps:fold}{向斜}。
(d) 按证明中的记号，正立翼的\omterm{def:b1:geological-maps:attitude}{倾角}为
$\gamma - \iota/2 = \ang{40}$，而倒转翼为 $\gamma + \iota/2 =
\ang{80}$，故 $\gamma = \ang{60}$，与\omterm{def:b1:faults-folds-fabrics:foliation}{劈理}所示相符，且 $\iota =
\ang{40}$：\omterm{def:b1:faults-folds-fabrics:fold-classes}{紧闭褶皱}。\omterm{def:b1:faults-folds-fabrics:fold-elements}{轴面}\omterm{def:b1:geological-maps:attitude}{走向}为北，\omterm{def:b1:geological-maps:attitude}{倾角} \ang{60}，\omterm{def:b1:geological-maps:attitude}{倾向}
为西，\omterm{def:b1:geological-maps:attitude}{产状} $270/60$。短翼已越过竖直方向，
其顶端朝东，\omterm{def:b1:geological-maps:fold}{褶皱}向东倾覆，与本章图中
所示一致。
\end{solution}

\begin{solution}{exo:b1:faults-folds-fabrics:12}
(a) 在对称的人字形\omterm{def:b1:geological-maps:fold}{褶皱}中，每一翼与
\omterm{def:b1:faults-folds-fabrics:fold-elements}{轴面}的夹角为 $\iota/2$，而\omterm{def:b1:faults-folds-fabrics:fold-elements}{轴面}垂直于缩短方向；长为 $l$ 的翼
因此横跨\omterm{def:b1:geological-maps:fold}{褶皱}的宽度为 $l\sin(\iota/2)$，整层则为
$L_0\sin(\iota/2)$。(b) $\sin\ang{30} = 0.5$：缩短量为
\qty{50}{\percent}。(c) $\sin(\iota/2) = 0.7$ 得 $\iota = \ang{88.9}$。
(d) $\iota = \ang{156}$ 且 $\sin\ang{78} = 0.978$：缩短量为
\qty{2.2}{\percent}，亦即 $1 - \cos\ang{12}$。
\end{solution}

\begin{solution}{pb:b1:faults-folds-fabrics:1}
\textbf{1.} $\mathbf n = (\sin\delta\sin\alpha, \sin\delta\cos\alpha,
\cos\delta)$ 可得
\[
  \mathbf n_1 = (0.20855, 0.03415, 0.97742), \qquad
  \mathbf n_2 = (-0.20855, 0.03415, 0.97742).
\]

\textbf{2.} $\mathbf n_1\cdot\mathbf n_2 = -0.04349 + 0.00117 + 0.95535 =
0.91302$，法向量之间的夹角为 \ang{24.1}。

\textbf{3.} 由\cref{prop:b1:faults-folds-fabrics:hinge}，
$\iota = \ang{180} - \ang{24.1} = \ang{155.9}$。

\textbf{4.} 在 \ang{170} 与 \ang{90} 之间：\omterm{def:b1:faults-folds-fabrics:fold-classes}{开阔褶皱}。

\textbf{5.} 两翼的\omterm{def:b1:geological-maps:attitude}{倾向}分别为 $080.7$ 与 $279.3$，并非正东和正西，
因为\omterm{def:b1:geological-maps:fold}{褶皱}倾伏：其最陡线不位于剖面内，剖面只
显示\omterm{def:b1:cross-sections:apparent-dip}{视倾角}。在东西向剖面中，由
\cref{thm:b1:cross-sections:apparent-dip}，$\tan\beta = \tan\ang{12.2}
\cos\ang{9.3} = 0.2134$，故 $\beta = \ang{12.0}$，且 $\ang{180} -
2\times\ang{12.0} = \ang{155.9}$（剖面相对该剖面倾斜 \ang{2}，
但在此精度下对结果无影响）。

\textbf{6.} 其两翼与\omterm{def:b1:geological-maps:fold}{背斜}的两翼平行，
故其\omterm{def:b1:faults-folds-fabrics:fold-elements}{翼间角}同样为 \ang{155.9}。

\textbf{7.} $\mathbf n_1\times\mathbf n_2 = (0.03415\times0.97742 -
0.97742\times0.03415,\; 0.97742\times(-0.20855) - 0.20855\times0.97742,\;
0.20855\times0.03415 + 0.03415\times0.20855) = (0, -0.40768, 0.01424)$.
取向下方向，$(0, 0.40768, -0.01424)$：\omterm{def:b1:faults-folds-fabrics:fold-elements}{枢纽}线\omterm{def:b1:geological-maps:attitude}{走向} $000$，
并向北倾伏，\omterm{def:b1:field-lab-methods-1:trend-plunge}{倾伏角} $\arctan(0.01424/0.40768) = \ang{2.0}$。

\textbf{8.} \omterm{def:b1:faults-folds-fabrics:fold-elements}{轴面}包含\omterm{def:b1:faults-folds-fabrics:fold-elements}{枢纽}线并平分两翼，
其法向沿 $\mathbf n_1 - \mathbf n_2 = (0.41710, 0, 0)$ 方向，水平
且为东西向。\omterm{def:b1:faults-folds-fabrics:fold-elements}{轴面}竖直，\omterm{def:b1:geological-maps:attitude}{走向}为南北：$090/90$。
该\omterm{def:b1:geological-maps:fold}{背斜}为\omterm{def:b1:faults-folds-fabrics:fold-classes}{直立褶皱}。

\textbf{9.} 开阔的\omterm{def:b1:faults-folds-fabrics:fold-classes}{直立褶皱}，向北倾伏 \ang{2.0}。

\textbf{10.} 由\cref{prop:b1:geological-maps:folds}可知，\omterm{def:b1:geological-maps:fold}{背斜}的\omterm{def:b1:geological-maps:geological-map}{露头}
在倾伏方向上闭合，即向北闭合。其核部保存最老的
岩石：覆盖层底部的砂岩，以及河流切割最深处
之下的基底片岩与花岗岩。

\textbf{11.} 该\omterm{def:b1:faults-folds-fabrics:foliation}{劈理}为\omterm{def:b1:faults-folds-fabrics:foliation}{轴面劈理}，竖直，\omterm{def:b1:geological-maps:attitude}{走向}南北，
$090/90$。由\cref{prop:b1:faults-folds-fabrics:intersection}，其
与层理的\omterm{def:b1:faults-folds-fabrics:lineation}{交线理}平行于\omterm{def:b1:faults-folds-fabrics:fold-elements}{枢纽}线：\omterm{def:b1:field-lab-methods-1:trend-plunge}{倾伏向}
$000$，\omterm{def:b1:field-lab-methods-1:trend-plunge}{倾伏角} \ang{2.0}。直接计算：平面 $080.7/12.2$ 与
\omterm{def:b1:geological-maps:attitude}{走向}北的竖直面相交于一条\omterm{def:b1:field-lab-methods-1:trend-plunge}{倾伏角}为 $\arctan(\tan\ang{12.2}
\cos\ang{80.7}) = \ang{2.0}$ 的线，\omterm{def:b1:field-lab-methods-1:trend-plunge}{倾伏向}为 $000$。

\textbf{12.} 面朝北看，东翼向左抬升 \ang{12.0}，
\omterm{def:b1:faults-folds-fabrics:foliation}{劈理}竖直：由\omterm{def:b1:faults-folds-fabrics:foliation}{劈理}到层理的锐角旋转，
角度为 \ang{78.0}，为逆时针方向，使\omterm{def:b1:faults-folds-fabrics:foliation}{劈理}顶端
向西移动。该翼为正立，故由
\cref{prop:b1:faults-folds-fabrics:cleavage-bedding}(b)，\omterm{def:b1:geological-maps:fold}{背斜}\omterm{def:b1:faults-folds-fabrics:fold-elements}{枢纽}
位于西侧 $x = \qty{3.5}{km}$ 处，与图中所示一致。

\textbf{13.} \omterm{def:b1:geological-maps:attitude}{走向}为 $120 - 90 = 030$；该面\omterm{def:b1:geological-maps:attitude}{倾向} $120$，故
\omterm{def:b1:faults-folds-fabrics:fault-elements}{上盘}位于迹线的东南东侧。

\textbf{14.} 该断层为\omterm{def:b1:faults-folds-fabrics:fault-types}{正断层}：$\sigma_1$ 竖直。

\textbf{15.} $\sigma_2$ 沿\omterm{def:b1:geological-maps:attitude}{走向}水平，即 $030$--$210$；
$\sigma_3$ 水平且垂直于\omterm{def:b1:geological-maps:attitude}{走向}，即 $120$--$300$。

\textbf{16.} 由 $\ang{45} + \phi/2 = \ang{60}$ 得 $\phi = \ang{30}$，
则 $\mu = \tan\ang{30} = 0.577$。

\textbf{17.} 共轭组同样\omterm{def:b1:geological-maps:attitude}{走向} $030$，\omterm{def:b1:geological-maps:attitude}{倾向}相反、\omterm{def:b1:geological-maps:attitude}{倾角}为 \ang{60}，
即 $300/60$。两组相交的夹角为 $\ang{180} - 2\times\ang{60} = \ang{60}$，
竖直的 $\sigma_1$ 平分该锐角。

\textbf{18.} 取 $\mu = 0.85$，$\phi = \ang{40.4}$，最有利取向的
\omterm{def:b1:faults-folds-fabrics:fault-types}{正断层}\omterm{def:b1:geological-maps:attitude}{倾角}为 \ang{65.2}；瓦尔登断层\omterm{def:b1:geological-maps:attitude}{倾角}为 \ang{60}，偏差 \ang{5.2}。
由\cref{exo:b1:faults-folds-fabrics:10}，其滑动时的应力比为
$\sigma_1'/\sigma_3' = (1 + 0.85\cot\ang{30})/(1 - 0.85\tan\ang{30}) = 4.85$，
而最有利面对应的值为 $4.68$：该断层接近但并非恰为最有利取向，
且仅需应力比高出 \qty{4}{\percent}，
便会再次滑动。

\textbf{19.} 滑动沿\omterm{def:b1:geological-maps:attitude}{倾向}向下：$r = \ang{-90}$；\omterm{def:b1:faults-folds-fabrics:fault-elements}{滑动矢量}的\omterm{def:b1:field-lab-methods-1:trend-plunge}{倾伏向}为
$120$，\omterm{def:b1:field-lab-methods-1:trend-plunge}{倾伏角}为 \ang{60}。

\textbf{20.} $(\cos\ang{60}\sin\ang{120}, \cos\ang{60}\cos\ang{120},
-\sin\ang{60}) = (0.4330, -0.2500, -0.8660)$。

\textbf{21.} $e = 200/\tan\ang{60} = \qty{115.5}{m}$。

\textbf{22.} 向东：$115.47\sin\ang{120} = \qty{100.0}{m}$；向北：
$115.47\cos\ang{120} = \qty{-57.7}{m}$，即向南 \qty{57.7}{m}。

\textbf{23.} $L = 200/\sin\ang{60} = \qty{230.9}{m}$，且
$200^2 + 115.47^2 = 40\,000 + 13\,333 = 53\,333 = 230.94^2$。\omterm{def:b1:cross-sections:throw}{垂直断距}与
\omterm{def:b1:cross-sections:throw}{水平断距}分别为\omterm{def:b1:faults-folds-fabrics:fault-elements}{滑动矢量}的竖直与水平分量，该矢量位于
\omterm{def:b1:geological-maps:attitude}{倾向}所在的竖直面内：二者是直角三角形的两条直角边，其
斜边即为滑动量。

\textbf{24.} 精确到米，瓦尔登断层的\omterm{def:b1:faults-folds-fabrics:fault-elements}{总滑距}为
\textbf{\qty[text-series-to-math=true]{231}{m}}。
\end{solution}
